%==============================================================================
% CAPÍTULO 25 — ÉTICA, LGPD E REGULAÇÃO DA IA EM ODONTOLOGIA
% PARTE V — Futuro, Ética e Regulação da IA em Odontologia
%==============================================================================
\chapter{Ética, LGPD e Regulação da Inteligência Artificial em Odontologia}
\label{cap:etica-lgpd-regulacao}

\nivelcinco

A inteligência artificial na Odontologia não é apenas uma questão técnica --
é, fundamentalmente, uma questão ética e jurídica. Cada radiografia panorâmica
analisada por uma rede neural, cada prontuário eletrônico processado por um
LLM, cada predição de risco de doença periodontal gerada por um modelo de
\textit{machine learning} carrega consigo um feixe de responsabilidades legais
que o cirurgião-dentista brasileiro precisa conhecer.

Este capítulo aborda o arcabouço regulatório que rege o uso de IA em saúde
bucal no Brasil, da Lei Geral de Proteção de Dados (LGPD) ao Projeto de Lei
2338/2023, passando pela regulação internacional (FDA, CE Mark) e pelos
princípios éticos que devem guiar cada profissional que opta por incorporar
sistemas inteligentes à sua prática clínica.

\indice{LGPD} \indice{ética em IA} \indice{regulação} \indice{PL 2338/2023}

%-----------------------------------------------------------------------------
% SEÇÃO 25.1
%-----------------------------------------------------------------------------
\section{LGPD e Dados de Saúde Bucal no Brasil}
\label{sec:lgpd-saude-bucal}

\nivelquatro

A Lei Geral de Proteção de Dados Pessoais (Lei nº 13.709, de 14 de agosto de
2018), em vigor desde setembro de 2020, estabelece que \textbf{dados de saúde
são dados pessoais sensíveis} (Art. 5º, inciso II) e, como tais, gozam de
proteção jurídica reforçada \cite{brasil2018lgpd}.

\indice{LGPD|(}

\subsection{O Que a LGPD Classifica como Dado de Saúde Bucal?}

Para a LGPD, são considerados dados de saúde bucal -- e, portanto, dados
pessoais sensíveis -- quaisquer informações relacionadas à saúde oral de uma
pessoa natural identificada ou identificável, incluindo, mas não se limitando a:

\begin{itemize}
  \item Radiografias (periapicais, panorâmicas, oclusais, TCFC) e imagens
    intraorais -- mesmo quando anonimizadas de forma reversível.
  \item Modelos digitais 3D de arcadas dentárias (escaneamento intraoral).
  \item Prontuários odontológicos (eletrônicos ou físicos) contendo diagnóstico,
    plano de tratamento, evolução clínica e prescrições.
  \item Imagens de lesões bucais, língua, mucosa -- inclusive fotografias
    clínicas com finalidade diagnóstica ou de documentação.
  \item Dados genéticos obtidos de amostras salivares com finalidade diagnóstica.
  \item Qualquer informação que, combinada com dados odontológicos, permita a
    identificação do titular (nome, CPF, CNS, endereço, telefone).
\end{itemize}

\destaque{A jurisprudência da ANPD (Autoridade Nacional de Proteção de Dados)
tem consolidado o entendimento de que imagens da face -- incluindo sorriso,
arcada dentária exposta e região peribucal -- constituem dado biométrico e,
portanto, dado pessoal sensível (Art. 5º, inciso V, LGPD). Radiografias
panorâmicas, que capturam toda a estrutura maxilofacial, inequivocamente se
enquadram nesta categoria.}

\subsection{Bases Legais para o Tratamento de Dados Odontológicos}

A LGPD estabelece que o tratamento de dados pessoais sensíveis só pode ocorrer
nas seguintes hipóteses (Art. 11):

\begin{enumerate}
  \item \textbf{Consentimento explícito e específico do titular} (Inciso I):
    O paciente deve ser informado, de forma clara e em linguagem acessível,
    sobre quais dados serão coletados, para quais finalidades, por quanto tempo
    serão armazenados e com quem poderão ser compartilhados. O consentimento
    deve ser obtido por escrito ou por meio eletrônico que demonstre a
    manifestação inequívoca de vontade.

  \item \textbf{Tutela da saúde} (Inciso II, alínea `f'): O tratamento de dados
    sensíveis é permitido sem consentimento quando for indispensável para a
    tutela da saúde, exclusivamente em procedimento realizado por profissionais
    de saúde, serviços de saúde ou autoridade sanitária. Esta base legal é a
    que fundamenta o tratamento de dados no contexto assistencial -- a consulta
    odontológica propriamente dita.

  \item \textbf{Pesquisa} (Inciso II, alínea `c'): A LGPD permite o tratamento
    de dados sensíveis para fins de pesquisa sem consentimento, desde que
    anonimizados. A anonimização, neste contexto, deve ser irreversível -- ou
    seja, não deve ser possível, por meios técnicos razoáveis, reidentificar o
    titular.

  \item \textbf{Exercício regular de direitos} (Inciso II, alínea `d'): Inclui
    o direito de defesa em processo judicial, administrativo ou arbitral,
    relevante para ações de responsabilidade civil.
\end{enumerate}

\notaexplicativa{A distinção entre anonimização (irreversível) e
pseudonimização (reversível, pois o dado pode ser reidentificado com uma chave)
é central na LGPD. Dados pseudonimizados continuam sendo dados pessoais e
estão sujeitos a todas as obrigações da lei. Apenas dados efetivamente
anonimizados -- sem possibilidade razoável de reidentificação -- deixam de ser
considerados dados pessoais (Art. 12).}

\subsection{Implicações Práticas para o Consultório que Usa IA}

\begin{table}[H]
  \centering
  \caption{Obrigações do cirurgião-dentista sob a LGPD ao utilizar IA.}
  \label{tab:lgpd-ia-consultorio}
  \begin{tabularx}{\textwidth}{lX}
    \toprule
    \textbf{Obrigação} & \textbf{Descrição} \\
    \midrule
    Consentimento & Obter consentimento explícito do paciente para uso de
      seus dados por sistemas de IA, especificando a finalidade \\
    Finalidade & Não utilizar dados coletados para diagnóstico em treinamento
      de modelos sem consentimento adicional \\
    Transparência & Informar ao paciente quando uma decisão clínica foi
      auxiliada por IA (Art. 20, LGPD) \\
    Segurança & Adotar medidas técnicas (criptografia, controle de acesso) e
      administrativas para proteger dados odontológicos \\
    Minimização & Coletar apenas os dados estritamente necessários para a
      finalidade declarada \\
    Eliminação & Eliminar dados após o término da finalidade ou do prazo
      legal de guarda (20 anos para prontuários, CFM/CFO) \\
    Direitos do titular & Garantir ao paciente acesso, correção,
      portabilidade e eliminação de seus dados \\
    \bottomrule
  \end{tabularx}
\end{table}

\subsection{Transferência Internacional de Dados Odontológicos}

Quando um cirurgião-dentista brasileiro utiliza serviços de IA hospedados em
nuvem no exterior (ex.: Google Colab, AWS, Azure, OpenAI API), está realizando
transferência internacional de dados pessoais sensíveis. A LGPD (Art. 33) exige
que o país de destino proporcione grau de proteção de dados adequado (decisão
de adequação da ANPD) ou que o controlador adote garantias contratuais
específicas (\textit{Standard Contractual Clauses}).

Na prática odontológica, isso significa:

\begin{itemize}
  \item Verificar se o fornecedor de IA possui servidores no Brasil (reduzindo
    a complexidade legal) ou, no mínimo, cláusulas contratuais compatíveis com
    a LGPD.
  \item Priorizar soluções de IA que processem dados \textit{on-premise} (no
    próprio consultório) ou em nuvem brasileira para minimizar riscos de
    transferência internacional.
  \item Documentar todas as transferências no Relatório de Impacto à Proteção
    de Dados Pessoais (RIPD), recomendado -- e em breve obrigatório -- pela
    ANPD para tratamento de dados sensíveis em larga escala.
\end{itemize}

\indice{LGPD|)}

%-----------------------------------------------------------------------------
% SEÇÃO 25.2
%-----------------------------------------------------------------------------
\section{Consentimento Informado para Inteligência Artificial}
\label{sec:consentimento-ia}

\nivelquatro

O consentimento informado -- pilar da bioética desde o Código de Nuremberg
(1947) e a Declaração de Helsinki (1964) -- adquire complexidade adicional
quando sistemas de IA estão envolvidos no cuidado odontológico.

\subsection{As Três Dimensões do Consentimento para IA}

\begin{enumerate}
  \item \textbf{Consentimento clínico (relação dentista-paciente):} O paciente
    consente com o uso de IA como ferramenta auxiliar no seu diagnóstico ou
    plano de tratamento. Deve abranger: (a) quais sistemas de IA serão
    utilizados, (b) qual o papel do sistema (auxiliar/suporte, não substituto
    do julgamento profissional), (c) limitações conhecidas do sistema
    (taxa de erro, populações para as quais o sistema não foi validado).

  \item \textbf{Consentimento para treinamento do modelo (pesquisa):} Quando os
    dados do paciente são utilizados para treinar ou aperfeiçoar um modelo de
    IA, o consentimento deve ser específico para esta finalidade. O paciente
    deve ser informado se seus dados serão compartilhados com terceiros
    (ex.: empresas de tecnologia, universidades parceiras).

  \item \textbf{Consentimento para armazenamento em nuvem (LGPD):} Se os dados
    do paciente serão processados ou armazenados em infraestrutura de nuvem --
    especialmente se hospedada fora do Brasil -- o paciente deve consentir com
    esta transferência internacional de dados sensíveis.
\end{enumerate}

\subsection{Elementos Obrigatórios do Termo de Consentimento para IA}

\begin{table}[H]
  \centering
  \caption{Elementos mínimos do Termo de Consentimento Livre e Esclarecido
  (TCLE) para uso de IA em Odontologia.}
  \label{tab:elementos-tcle-ia}
  \begin{tabularx}{\textwidth}{lX}
    \toprule
    \textbf{Elemento} & \textbf{Conteúdo Mínimo} \\
    \midrule
    Identificação do sistema & Nome, versão, desenvolvedor e registro
      (ANVISA, se aplicável) \\
    Finalidade específica & Qual tarefa a IA executa (detecção, classificação,
      predição, sugestão de plano) \\
    Papel do sistema & Auxiliar/suporte -- NÃO substitui o julgamento do
      profissional \\
    Dados coletados & Lista exaustiva de quais dados serão processados
      (imagens, texto, áudio) \\
    Destino dos dados & Onde serão armazenados (servidor local, nuvem, país) \\
    Compartilhamento & Com quem os dados poderão ser compartilhados \\
    Prazo de armazenamento & Por quanto tempo os dados serão mantidos \\
    Direitos do titular & Acesso, correção, eliminação, revogação do
      consentimento \\
    Riscos conhecidos & Vieses do modelo, taxa de erro reportada, limitações
      populacionais \\
    Contato do responsável & Nome, CRM/CRO, telefone e e-mail do controlador \\
    \bottomrule
  \end{tabularx}
\end{table}

\citacaoativa{doi-1016-j-adaj-2023-06-002}{10.1016/j.adaj.2023.06.002}{%
  ``O consentimento informado para IA odontológica não pode ser genérico. O
  paciente tem o direito de saber, em linguagem compreensível, que um algoritmo
  -- e não apenas o julgamento humano do dentista -- contribuiu para a decisão
  que afeta sua saúde bucal.''}

%-----------------------------------------------------------------------------
% SEÇÃO 25.3
%-----------------------------------------------------------------------------
\section{Anonimização e Pseudonimização de Imagens Odontológicas}
\label{sec:anonimizacao}

\nivelcinco

A anonimização de imagens odontológicas para pesquisa e desenvolvimento de IA
é um desafio técnico e jurídico de primeira ordem. Ao contrário do que sugere o
senso comum, remover o nome do arquivo DICOM ou cortar metadados não é
suficiente para anonimizar uma imagem odontológica.

\indice{anonimização} \indice{pseudonimização}

\subsection{Por que Imagens Odontológicas São Difíceis de Anonimizar?}

Uma radiografia panorâmica contém, intrinsecamente, informações que podem
permitir a reidentificação do paciente:

\begin{itemize}
  \item \textbf{Anatomia única:} A configuração de dentes presentes, ausentes,
    restaurados e implantados -- a chamada ``impressão digital odontológica'' --
    é virtualmente única para cada indivíduo, análoga a uma impressão digital
    datiloscópica.
  \item \textbf{Materiais dentários identificáveis:} Coroas, pontes, implantes
    e restaurações com formatos e posições específicas constituem assinaturas
    individuais.
  \item \textbf{Patologias e anomalias raras:} A combinação de condições
    odontológicas incomuns pode ser suficientemente distintiva para
    reidentificação.
  \item \textbf{Metadados DICOM:} Arquivos DICOM contêm dezenas de campos com
    informações potencialmente identificadoras (nome do paciente, data de
    nascimento, ID do paciente, instituição, data do exame, parâmetros técnicos
    do equipamento).
  \item \textbf{Sobreposição facial:} Radiografias panorâmicas e TCFC capturam
    estruturas ósseas da face que, combinadas com outras fontes de dados,
    podem permitir reidentificação.
\end{itemize}

\destaque{A ANPD e o \textit{European Data Protection Board} (EDPB) convergem no
entendimento de que a mera remoção de identificadores diretos (nome, CPF) NÃO
constitui anonimização efetiva se existirem meios técnicos razoáveis de
reidentificação. Imagens médicas/odontológicas, por sua natureza rica em
informações, raramente podem ser consideradas plenamente anonimizadas.}

\subsection{Estratégias Técnicas de Desidentificação}

\begin{description}
  \item[Limpeza de metadados DICOM:] Remoção ou substituição de todos os
    \textit{tags} DICOM que contenham informações do paciente (PatientName,
    PatientID, PatientBirthDate, InstitutionName, StudyDate, etc.). Ferramentas
    como \texttt{pydicom} (Python) e \texttt{dcm4che} (Java) oferecem funções
    específicas para anonimização DICOM em lote.

  \item[Desfoque seletivo (\textit{selective blurring}):] Aplicação de
    desfoque gaussiano ou \textit{pixelation} em regiões potencialmente
    identificadoras, como a região do seio frontal e estruturas ósseas faciais
    não relevantes para a tarefa clínica. Esta técnica deve ser aplicada com
    cautela para não degradar as regiões de interesse odontológico.

  \item[Máscara facial (\textit{face masking}):] Em TCFC, algoritmos de
    segmentação podem isolar e remover a superfície facial, preservando apenas
    as estruturas dentoalveolares.

  \item[Geração sintética:] Uso de \textit{Generative Adversarial Networks}
    (GANs) ou \textit{Diffusion Models} para gerar imagens odontológicas
    sintéticas que preservam as propriedades estatísticas do \textit{dataset}
    original sem corresponder a nenhum paciente real. O OdontoIA mantém um
    \textit{pipeline} aberto de geração de radiografias sintéticas.

  \item[K-anonimato em \textit{datasets} odontológicos:] Garantir que cada
    combinação de características quase-identificadoras (faixa etária, sexo,
    região geográfica, características odontológicas raras) apareça em pelo
    menos $k$ registros distintos no \textit{dataset}.
\end{description}

\notaexplicativa{O conceito de \textit{k-anonimato}, introduzido por Sweeney
(2002), estabelece que um \textit{dataset} é \textit{k}-anônimo se cada registro
é indistinguível de pelo menos $k-1$ outros registros com respeito a um conjunto
de quase-identificadores. Para \textit{datasets} odontológicos, recomendamos
$k \geq 5$.}

\subsection{O Papel do OdontoIA na Desidentificação}

O ecossistema OdontoIA oferece, como parte de seu compromisso com a ética e a
conformidade legal, um módulo de desidentificação de imagens odontológicas que:

\begin{itemize}
  \item Remove automaticamente metadados DICOM sensíveis (\texttt{pydicom}).
  \item Aplica máscara facial em TCFC (modelo U-Net treinado em 500 exames).
  \item Gera relatório de auditoria de anonimização exportável.
  \item Sugere, quando a anonimização completa não é possível, a via da
    pseudonimização com Termo de Compromisso de Utilização de Dados (TCUD).
\end{itemize}

%-----------------------------------------------------------------------------
% SEÇÃO 25.4
%-----------------------------------------------------------------------------
\section{Viés Algorítmico: Etnia, Gênero, Idade e Classe Social}
\label{sec:vies-algoritmico}

\nivelcinco

O viés algorítmico é, talvez, o desafio ético mais insidioso da IA em saúde --
porque é invisível a olho nu e, frequentemente, se disfarça de objetividade
matemática. Quando um modelo de \textit{deep learning} apresenta acurácia de
93\% na detecção de lesões bucais, a pergunta que precisa ser feita é:
93\% para quem?

\indice{viés algorítmico|(}

\subsection{As Quatro Fontes de Viés em IA Odontológica}

\begin{enumerate}
  \item \textbf{Viés de representação (\textit{representation bias}):} Ocorre
    quando o \textit{dataset} de treinamento não reflete a diversidade da
    população-alvo. Exemplo concreto: um modelo treinado exclusivamente com
    radiografias de pacientes caucasianos europeus terá desempenho
    significativamente inferior em pacientes negros, asiáticos e indígenas
    brasileiros, cujas características de densidade óssea, anatomia
    dentofacial e padrões de doença podem diferir.

  \item \textbf{Viés de anotação (\textit{label bias}):} Ocorre quando o
    \textit{ground truth} é estabelecido por anotadores que carregam seus
    próprios vieses. Se um \textit{dataset} de lesões orais é anotado apenas por
    especialistas formados em uma única escola odontológica, os critérios
    diagnósticos podem não ser representativos da diversidade de práticas
    clínicas.

  \item \textbf{Viés histórico (\textit{historical bias}):} Ocorre quando os
    dados refletem desigualdades estruturais pré-existentes. Se um modelo de
    predição de risco de cárie é treinado com dados de pacientes que tiveram
    acesso regular a tratamento odontológico, ele sistematicamente subestimará
    o risco de pacientes que não tiveram esse acesso -- tipicamente, populações
    de baixa renda atendidas pelo SUS.

  \item \textbf{Viés de confirmação (\textit{confirmation bias}):} Ocorre quando
    o clínico confia excessivamente nas predições do modelo, ignorando
    evidências contraditórias do exame clínico. Este viés é particularmente
    perigoso em sistemas com interface ``caixa-preta'' que não oferecem
    explicações para suas predições.
\end{enumerate}

\subsection{Evidências de Viés em IA Odontológica}

Um estudo de 2023 analisou 45 modelos de \textit{deep learning} para
diagnóstico odontológico publicados entre 2018--2022 e constatou que:

\begin{itemize}
  \item \textbf{87\%} dos \textit{datasets} não reportaram a composição
    étnico-racial da amostra.
  \item \textbf{92\%} não realizaram análise de desempenho por subgrupo (sexo,
    idade, etnia).
  \item \textbf{76\%} utilizaram dados de um único país -- predominantemente
    EUA (34\%), China (22\%) e Coreia do Sul (12\%).
  \item \textbf{Apenas 2 estudos} (4,4\%) incluíram pacientes brasileiros.
  \item \textbf{Nenhum estudo} reportou nível socioeconômico dos pacientes ou
    tipo de cobertura de saúde (pública, privada, filantrópica).
\end{itemize}

Esses números são alarmantes. Eles significam que, em 2024, a vasta maioria dos
modelos de IA odontológica publicados não pode ter seu desempenho generalizado
para a população brasileira -- muito menos para as populações historicamente
marginalizadas que mais necessitam de melhorias no acesso ao diagnóstico.

\citacaoativa{doi-1038-s41591-021-01614-0}{10.1038/s41591-022-01981-2}{%
  ``Modelos de IA treinados em populações homogêneas podem exacerbar
  desigualdades em saúde quando implantados em populações diversas. A auditoria
  de equidade deve ser um requisito regulatório, não uma opção ética.''}

\subsection{Estratégias de Mitigação de Viés}

\begin{description}
  \item[Auditoria de equidade:] Antes de publicar ou implantar um modelo,
    calcule métricas de desempenho estratificadas por sexo, faixa etária,
    etnia/raça (categorias IBGE), nível socioeconômico e tipo de serviço de
    saúde (público/privado). Se houver diferença estatisticamente significativa
    ($p < 0,05$) entre subgrupos, o modelo não está pronto para implantação.

  \item[Coleta intencionalmente diversa:] Projete \textit{datasets} que reflitam
    a diversidade da população brasileira. O Censo IBGE 2022 reporta 45,3\% de
    pardos, 43,5\% de brancos, 10,2\% de pretos, 0,6\% de indígenas e 0,4\% de
    amarelos. Seu \textit{dataset} se aproxima dessa distribuição?

  \item[\textit{Data augmentation} consciente:] Utilize técnicas de aumento de
    dados que compensem sub-representação de grupos minoritários, mas sem
    introduzir artefatos irreais.

  \item[\textit{Fairness constraints} no treinamento:] Incorpore restrições de
    equidade diretamente na função de perda do modelo, penalizando
    discrepâncias de desempenho entre subgrupos.

  \item[Transparência radical:] Publique a \textit{datasheet} completa do
    \textit{dataset} (seguindo o modelo de \cite{gebru2021datasheets}),
    incluindo composição demográfica, critérios de inclusão/exclusão e vieses
    conhecidos.
\end{description}

\indice{viés algorítmico|)}

%-----------------------------------------------------------------------------
% SEÇÃO 25.5
%-----------------------------------------------------------------------------
\section{Regulação da IA no Brasil: PL 2338/2023}
\label{sec:pl-2338}

\nivelquatro

O Projeto de Lei nº 2.338/2023, em tramitação no Senado Federal, representa a
tentativa mais abrangente do legislador brasileiro de estabelecer um marco
regulatório para a inteligência artificial. Inspirado no AI Act europeu, mas
com adaptações à realidade brasileira, o PL 2338 terá impacto direto sobre o
uso de IA em Odontologia.

\indice{PL 2338/2023|(}

\subsection{Estrutura do PL 2338/2023}

O projeto estrutura-se em torno de três eixos principais:

\begin{enumerate}
  \item \textbf{Classificação de riscos:} Os sistemas de IA são classificados
    em níveis de risco -- risco excessivo (proibidos), alto risco (regulação
    intensa), risco limitado (obrigações de transparência) e risco mínimo (sem
    obrigações adicionais). Sistemas de IA para diagnóstico e tratamento em
    saúde são explicitamente categorizados como \textbf{alto risco} (Art. 17,
    inciso I).

  \item \textbf{Direitos das pessoas afetadas:} Direito à explicação (Art. 8º),
    direito à revisão humana de decisões automatizadas (Art. 9º), direito à
    não-discriminação algorítmica (Art. 11) e direito à correção de dados
    incorretos (Art. 10).

  \item \textbf{Obrigações dos desenvolvedores e operadores:} Avaliação de
    impacto algorítmico (Art. 23), documentação técnica completa (Art. 24),
    supervisão humana (Art. 25), robustez e acurácia (Art. 26), transparência
    (Art. 27) e governança de dados (Art. 28).
\end{enumerate}

\subsection{Implicações para o Cirurgião-Dentista}

Se aprovado nos termos atuais, o PL 2338/2023 imporá ao cirurgião-dentista que
utiliza IA as seguintes obrigações:

\begin{table}[H]
  \centering
  \caption{Obrigações previstas no PL 2338/2023 para uso de IA em Odontologia.}
  \label{tab:pl2338-obrigacoes}
  \begin{tabularx}{\textwidth}{lX}
    \toprule
    \textbf{Obrigação} & \textbf{Descrição} \\
    \midrule
    Avaliação de impacto & Realizar avaliação de impacto algorítmico antes de
      implantar qualquer sistema de IA de alto risco \\
    Supervisão humana & Garantir que decisões automatizadas possam ser revistas
      por um profissional humano qualificado \\
    Transparência & Informar ao paciente, de forma clara e acessível, que um
      sistema de IA foi utilizado \\
    Documentação & Manter registro das versões do sistema, \textit{datasets} de
      treinamento, métricas de desempenho e auditorias \\
    Não-discriminação & Assegurar que o sistema não produza resultados
      discriminatórios com base em raça, gênero, idade ou classe social \\
    Responsabilização & Responder civilmente por danos causados por decisões
      baseadas em IA (Art. 6º, CDC, solidariedade da cadeia de fornecimento) \\
    \bottomrule
  \end{tabularx}
\end{table}

\subsection{Comparação com o AI Act Europeu}

O PL 2338/2023 brasileiro e o AI Act europeu (Regulamento UE 2024/1689)
compartilham a abordagem baseada em risco, mas divergem em aspectos importantes:

\begin{itemize}
  \item \textbf{Abrangência:} O AI Act aplica-se a \textit{providers}
    (desenvolvedores) e \textit{deployers} (usuários profissionais) de sistemas
    de IA. O PL 2338 adota terminologia similar: ``agentes de IA'' --
    desenvolvedores e operadores.

  \item \textbf{Classificação de dispositivos médicos:} Na União Europeia,
    \textit{software} com finalidade médica é regulado pelo MDR (\textit{Medical
    Device Regulation} 2017/745). O AI Act adiciona camada regulatória
    específica para IA, mas mantém o MDR como base. No Brasil, a ANVISA regula
    \textit{software} como dispositivo médico (RDC 657/2022), e o PL 2338
    adicionará a camada de IA.

  \item \textbf{\textit{Sandbox} regulatório:} Ambos preveem ambientes
    controlados de experimentação (\textit{regulatory sandboxes}). Este
    instrumento será crucial para que \textit{startups} odontológicas
    brasileiras como o OdontoIA possam testar inovações em ambiente
    supervisionado antes da certificação completa.
\end{itemize}

\destaque{O PL 2338/2023, em seu Art. 17, §2º, estabelece que a autoridade
competente definirá os critérios para classificação de sistemas de IA de alto
risco. Até a publicação desta regulamentação infralegal, recomenda-se que todo
sistema de IA com finalidade diagnóstica, prognóstica ou terapêutica em
Odontologia seja tratado \textit{como se já fosse} classificado como alto risco.}

\indice{PL 2338/2023|)}

%-----------------------------------------------------------------------------
% SEÇÃO 25.6
%-----------------------------------------------------------------------------
\section{FDA e CE Mark para Software Odontológico com IA}
\label{sec:fda-ce-mark}

\nivelcinco

Para o pesquisador ou empreendedor odontológico brasileiro que almeja
mercados internacionais, compreender os processos regulatórios do FDA (EUA) e
CE Mark (União Europeia) é essencial. Estes são, atualmente, os dois principais
selos de conformidade para \textit{software} médico com componente de IA.

\indice{FDA} \indice{CE Mark}

\subsection{FDA: \textit{Software as a Medical Device} (SaMD)}

O \textit{Center for Devices and Radiological Health} (CDRH) do FDA classifica
\textit{software} odontológico com IA como \textit{Software as a Medical Device}
(SaMD). O caminho regulatório depende da classe de risco:

\begin{description}
  \item[Classe I (risco baixo):] \textit{Software} de gerenciamento de imagens
    sem análise diagnóstica, \textit{software} educacional. Isento de
    \textit{Premarket Notification} (510k). A maioria dos aplicativos
    odontológicos de uso geral se enquadra aqui.

  \item[Classe II (risco moderado):] \textit{Software} de detecção auxiliada
    por computador (CADe/CADx), \textit{software} de análise de imagens
    radiográficas. Requer \textit{Premarket Notification} 510(k) demonstrando
    equivalência substancial a um dispositivo já legalmente comercializado
    (\textit{predicate device}).

  \item[Classe III (alto risco):] \textit{Software} que fornece diagnóstico
    definitivo sem revisão humana, \textit{software} que controla dispositivos
    implantáveis. Requer \textit{Premarket Approval} (PMA) com evidência
    clínica robusta.
\end{description}

A maioria dos sistemas de IA odontológica -- detecção de cáries, lesões
periapicais, análise ortodôntica -- enquadra-se na Classe II e requer 510(k).

\notaexplicativa{Em 2024, o FDA introduziu o conceito de \textit{Predetermined
Change Control Plan} (PCCP) para dispositivos com IA que continuam aprendendo
após a comercialização (\textit{adaptive AI/ML}). O PCCP permite que o
fabricante especifique, \textit{a priori}, os tipos de modificações que o
algoritmo poderá sofrer sem necessidade de nova submissão regulatória.}

\subsection{CE Mark: MDR 2017/745 e AI Act}

Na União Europeia, o caminho regulatório é duplo:

\begin{enumerate}
  \item \textbf{MDR (\textit{Medical Device Regulation} 2017/745):} Classifica
    o \textit{software} odontológico como dispositivo médico (Classe I, IIa, IIb
    ou III) com base no risco. \textit{Software} diagnóstico é tipicamente
    Classe IIa.

  \item \textbf{AI Act (Regulamento UE 2024/1689):} Adiciona requisitos
    específicos para sistemas de IA de alto risco, incluindo: qualidade dos
    \textit{datasets} de treinamento (Art. 10), documentação técnica (Art. 11),
    transparência (Art. 13), supervisão humana (Art. 14) e acurácia/robustez
    (Art. 15).
\end{enumerate}

\subsection{OdontoIA e o Caminho Regulatório Internacional}

O ecossistema OdontoIA foi concebido desde sua arquitetura inicial com vistas à
conformidade regulatória internacional. O módulo de rastreabilidade SDD/TDD
(Capítulos~12--17) gera automaticamente a documentação exigida por FDA 510(k) e
CE Mark, incluindo:

\begin{itemize}
  \item Especificação de requisitos de \textit{software} (SRS) em formato
    IEC 62304.
  \item Matriz de rastreabilidade entre requisitos, código e testes.
  \item Relatório de verificação e validação (V\&V) com cobertura de testes
    superior a 90\%.
  \item \textit{Datasheet} do \textit{dataset} de treinamento conforme modelo
    \textit{Datasheets for Datasets}.
  \item Plano de controle de mudanças predeterminado (PCCP) para modelos
    adaptativos.
\end{itemize}

%-----------------------------------------------------------------------------
% SEÇÃO 25.7
%-----------------------------------------------------------------------------
\section{Responsabilidade Civil do Cirurgião-Dentista que Usa IA}
\label{sec:responsabilidade-civil}

\nivelquatro

A pergunta que tira o sono de muitos profissionais: \textbf{se a IA errar o
diagnóstico e o paciente sofrer dano, de quem é a culpa?}

A resposta jurídica no ordenamento brasileiro, à luz do Código Civil (Lei nº
10.406/2002), do Código de Defesa do Consumidor (Lei nº 8.078/1990) e do
Código de Ética Odontológica (Resolução CFO 118/2012), é inequívoca:
\textbf{a responsabilidade é do cirurgião-dentista.}

\indice{responsabilidade civil|(}

\subsection{Fundamentos Jurídicos}

\begin{enumerate}
  \item \textbf{Obrigação de meio (não de resultado):} A Odontologia, como
    regra geral, é uma obrigação de meio -- o profissional se compromete a
    empregar os melhores esforços, técnica e ciência disponíveis, mas não pode
    garantir um resultado específico. A IA não altera essa natureza jurídica da
    relação dentista-paciente.

  \item \textbf{Responsabilidade subjetiva (com culpa):} O art. 951 do Código
    Civil e o art. 14, §4º do CDC estabelecem que a responsabilidade do
    profissional liberal é apurada mediante verificação de culpa (imperícia,
    imprudência ou negligência). O uso de IA não transforma a responsabilidade
    subjetiva em objetiva -- o paciente ainda precisa demonstrar que houve
    culpa do profissional.

  \item \textbf{Dever de atualização científica:} O Código de Ética Odontológica
    (Art. 5º, inciso IV) impõe ao cirurgião-dentista o dever de ``zelar pela
    saúde e dignidade do paciente'' e (Art. 6º, inciso III) ``manter-se
    atualizado com a evolução da ciência odontológica''. Utilizar IA
    desatualizada, não validada ou inadequada pode configurar imperícia.

  \item \textbf{\textit{Standard of care}:} O parâmetro de conduta profissional
    (\textit{standard of care}) evolui com a tecnologia. Se, em 2028, a
    detecção de lesões periapicais com IA for padrão na comunidade odontológica,
    o profissional que optar por não utilizar IA pode, paradoxalmente, estar
    abaixo do \textit{standard of care} -- desde que a tecnologia tenha
    evidência robusta de superioridade clínica.
\end{enumerate}

\subsection{Cenários de Responsabilidade}

\begin{table}[H]
  \centering
  \caption{Cenários de responsabilidade civil no uso de IA odontológica.}
  \label{tab:cenarios-responsabilidade}
  \begin{tabularx}{\textwidth}{p{3cm}XX}
    \toprule
    \textbf{Cenário} & \textbf{Responsabilidade} & \textbf{Fundamento} \\
    \midrule
    IA sugere diagnóstico correto, dentista ignora e erra &
      Dentista -- negligência &
      O profissional é o decisor final; ignorar evidência disponível
      configura culpa \\
    IA erra, dentista segue cegamente, paciente sofre dano &
      Dentista -- imperícia &
      Confiar em ferramenta sem verificação crítica viola o
      \textit{standard of care} \\
    IA erra por viés de \textit{dataset}, dentista não tinha como saber &
      Dentista (com possível ação regressiva contra fabricante) &
      Responsabilidade primária é do profissional; este pode buscar
      ressarcimento do fabricante \\
    IA funcionava corretamente na compra, fabricante atualiza e introduz erro &
      Fabricante (CDC, fato do produto) &
      Vício do produto por atualização defeituosa \\
    IA é usada \textit{off-label} (fora da finalidade aprovada) &
      Dentista -- imprudência &
      Uso de dispositivo fora das especificações do fabricante \\
    IA não tem registro ANVISA quando exigível &
      Dentista -- ilegalidade &
      Comercialização/uso de produto não registrado (Lei 6.360/76, Art. 12) \\
    \bottomrule
  \end{tabularx}
\end{table}

\subsection{Recomendações Práticas para Mitigação de Risco}

\begin{enumerate}
  \item \textbf{Documente o uso da IA no prontuário:} Registre qual sistema foi
    utilizado, qual versão, qual a predição fornecida e qual sua decisão final
    (concordante ou divergente da IA). Se divergir, justifique clinicamente.

  \item \textbf{Utilize apenas sistemas com registro sanitário}: Verifique se
    o \textit{software} possui registro ANVISA (RDC 657/2022) quando exigível.
    \textit{Software} exclusivamente educacional ou de gestão pode ser isento,
    mas \textit{software} diagnóstico/decisório não.

  \item \textbf{Mantenha seguro de responsabilidade civil profissional}: Verifique
    se sua apólice cobre danos decorrentes do uso de tecnologia de IA.

  \item \textbf{Não delegue a decisão final à IA}: O Conselho Federal de
    Odontologia (CFO) é inequívoco: o ato diagnóstico é privativo do
    cirurgião-dentista. A IA é ferramenta auxiliar, não substituta.

  \item \textbf{Atualize-se continuamente}: Participe de cursos, Congressos e
    comunidades como o OdontoIA. A evolução tecnológica é veloz, e o
    \textit{standard of care} evoluirá com ela.
\end{enumerate}

\citacaoativa{arnet2022cancer}{10.1016/j.oooo.2022.01.005}{%
  ``O dentista permanece, legal e eticamente, o responsável final pelo
  diagnóstico e plano de tratamento, independentemente do grau de automação
  da ferramenta que utiliza. A IA é um instrumento -- como o explorador, a
  sonda periodontal ou a radiografia -- e, como todo instrumento, seu valor
  depende da competência de quem o maneja.''}

\indice{responsabilidade civil|)}

%-----------------------------------------------------------------------------
% SEÇÃO 25.8
%-----------------------------------------------------------------------------
\section{Declaração de Helsinki e Pesquisa com IA}
\label{sec:helsinki-ia}

\nivelquatro

A Declaração de Helsinki (1964, última revisão 2024), da Associação Médica
Mundial (WMA), estabelece os princípios éticos para pesquisa envolvendo seres
humanos. Embora originalmente concebida para a medicina, seus princípios
aplicam-se integralmente à pesquisa odontológica -- e, por extensão, à pesquisa
com IA em Odontologia.

\indice{Declaração de Helsinki|(}

\subsection{Princípios Relevantes para Pesquisa com IA}

\begin{enumerate}
  \item \textbf{Primazia do bem-estar do participante} (Art. 4): Os interesses
    da ciência e da sociedade nunca devem prevalecer sobre o bem-estar do
    participante da pesquisa. Isso significa que, se houver conflito entre
    ``melhorar a acurácia do modelo'' e ``proteger a privacidade do
    participante'', a privacidade prevalece.

  \item \textbf{Aprovação por comitê de ética} (Art. 23): Toda pesquisa com
    seres humanos -- incluindo pesquisa retrospectiva com dados de prontuários
    e imagens -- deve ser submetida e aprovada por um comitê de ética em
    pesquisa (CEP) antes de seu início. No Brasil, a submissão é feita pela
    Plataforma Brasil (\url{https://plataformabrasil.saude.gov.br}).

  \item \textbf{Consentimento informado} (Art. 25--26): Os participantes devem
    ser informados sobre os objetivos, métodos, riscos e benefícios da pesquisa,
    incluindo o fato de que seus dados serão utilizados para treinar modelos de
    IA. O consentimento para uso futuro de dados (ex.: repositórios públicos)
    deve ser específico e destacado.

  \item \textbf{Privacidade e confidencialidade} (Art. 24): Medidas adequadas
    devem ser tomadas para proteger a privacidade dos participantes e a
    confidencialidade de seus dados -- um desafio particular em IA, onde a
    anonimização completa é difícil (conforme discutido na Seção~25.3).

  \item \textbf{Publicação de resultados} (Art. 36): Pesquisadores têm o dever
    ético de publicar seus resultados, sejam eles positivos, negativos ou
    inconclusivos. Modelos de IA que ``não funcionaram'' também geram
    conhecimento valioso e devem ser reportados -- o viés de publicação em IA
    é um problema científico e ético.
\end{enumerate}

\subsection{Pesquisa com Dados Retrospectivos: A Nova Fronteira Ética}

A maior parte da pesquisa em IA odontológica utiliza dados retrospectivos --
radiografias, prontuários e imagens coletados durante a assistência clínica.
Esta modalidade apresenta desafios éticos específicos:

\begin{itemize}
  \item \textbf{Dispensa de TCLE:} A Resolução CNS 466/2012 (Art. IV.8) permite
    a dispensa do Termo de Consentimento Livre e Esclarecido em pesquisas
    retrospectivas com dados de prontuário quando a obtenção do consentimento
    for logisticamente impraticável e a pesquisa não envolver riscos adicionais
    aos participantes. No entanto, a dispensa deve ser expressamente autorizada
    pelo CEP.

  \item \textbf{Compromisso de utilização de dados:} Em substituição ao TCLE,
    os pesquisadores devem assinar um Termo de Compromisso de Utilização de
    Dados (TCUD), comprometendo-se a:
    \begin{itemize}
      \item Utilizar os dados exclusivamente para a finalidade declarada.
      \item Não divulgar ou compartilhar dados que permitam identificação.
      \item Destruir os dados após o término da pesquisa (ou depositá-los
        anonimizados em repositório público, se autorizado pelo CEP).
      \item Reportar imediatamente ao CEP qualquer violação ou incidente de
        segurança.
    \end{itemize}
\end{itemize}

\subsection{OdontoIA e o Compromisso com a Ética em Pesquisa}

O ecossistema OdontoIA adota, como política institucional, os seguintes
compromissos éticos:

\begin{itemize}
  \item Todos os \textit{datasets} disponibilizados publicamente pelo OdontoIA
    são acompanhados do número do parecer CEP que autorizou sua coleta e
    compartilhamento.
  \item Nenhum \textit{dataset} do OdontoIA contém dados de pacientes que não
    tenham consentido (ou cujo TCLE não tenha sido dispensado formalmente pelo
    CEP).
  \item O OdontoIA não utiliza dados de pacientes para treinamento de seus
    modelos sem aprovação ética específica para esta finalidade.
  \item O módulo de auditoria ética do OdontoIA permite que pesquisadores gerem
    automaticamente a documentação necessária para submissão à Plataforma Brasil.
\end{itemize}

\indice{Declaração de Helsinki|)}

%-----------------------------------------------------------------------------
% SEÇÃO 25.9
%-----------------------------------------------------------------------------
\section{Prática: Template de Consentimento LGPD para Datasets Odontológicos}
\label{sec:pratica-lgpd}

\nivelcinco

\begin{pratica}{Criação de Termo de Consentimento LGPD para Dataset Odontológico}
  Objetivo: Elaborar um Termo de Consentimento Livre e Esclarecido (TCLE) que
  atenda simultaneamente aos requisitos da LGPD (Lei 13.709/2018), da Resolução
  CNS 466/2012 e do CLAIM 2024 para um \textit{dataset} de radiografias
  panorâmicas destinado ao treinamento de modelo de detecção de lesões ósseas.
\end{pratica}

\subsection{Template de TCLE -- Uso de Imagens Odontológicas para Treinamento de IA}

\begin{quote}
\textbf{TERMO DE CONSENTIMENTO LIVRE E ESCLARECIDO (TCLE)}

\textbf{Título da pesquisa:} Desenvolvimento de Modelo de Inteligência Artificial
para Detecção de Lesões Ósseas em Radiografias Panorâmicas

\textbf{Pesquisador responsável:} [Nome completo, CRO, instituição, e-mail,
telefone]

\textbf{Instituição:} [Nome da universidade/hospital/clínica]

\textbf{CAAE:} [Número do Certificado de Apresentação de Apreciação Ética]

\bigskip

\textbf{1. OBJETIVO DA PESQUISA}

Convidamos o(a) senhor(a) a autorizar o uso de sua radiografia panorâmica,
previamente realizada por indicação clínica, para o treinamento de um programa
de computador (inteligência artificial) que auxiliará cirurgiões-dentistas na
detecção de lesões ósseas nos maxilares.

\textbf{2. PROCEDIMENTOS}

A sua participação consiste unicamente em permitir que sua radiografia
panorâmica -- já existente em seu prontuário, realizada por necessidade clínica
-- seja incluída em um banco de imagens para treinamento do programa de
computador. NÃO haverá qualquer exame, procedimento ou consulta adicional.

\textbf{3. DADOS COLETADOS E FINALIDADE (LGPD, Art. 9º, §2º)}

Serão coletados: (a) sua radiografia panorâmica digital; (b) seu sexo (masculino
/ feminino); (c) sua idade em anos completos; (d) o laudo radiográfico emitido
pelo cirurgião-dentista que o(a) atendeu. Esses dados serão utilizados
EXCLUSIVAMENTE para treinar e testar o programa de computador.

\textbf{Seus dados NÃO serão utilizados para:} (i) qualquer outra finalidade que
não a descrita acima; (ii) tomada de decisão automatizada que afete seus
direitos; (iii) definição de perfil pessoal, profissional ou de crédito.

\textbf{4. ARMAZENAMENTO E COMPARTILHAMENTO (LGPD, Art. 9º, §4º)}

Seus dados serão armazenados em servidor seguro na [instituição/no Brasil],
protegidos por criptografia e controle de acesso restrito aos pesquisadores
autorizados. Seus dados PODERÃO ser compartilhados, de forma anonimizada, em
repositório científico público (ex.: Zenodo) para fins de reprodutibilidade da
pesquisa. Nenhum dado que permita sua identificação (nome, CPF, CNS, data de
nascimento completa) será compartilhado.

\textbf{5. ANONIMIZAÇÃO E PSEUDONIMIZAÇÃO (LGPD, Art. 12 e 13)}

Sua radiografia terá seu nome e demais identificadores removidos (metadados
DICOM), sendo substituídos por um código numérico aleatório. A chave que liga
seu nome ao código será armazenada separadamente e destruída após 5 anos. Embora
medidas técnicas rigorosas sejam adotadas, é importante que o(a) senhor(a) saiba
que a anonimização completa de imagens radiográficas é tecnicamente desafiadora
-- a configuração única de seus dentes e estruturas ósseas pode, em teoria,
permitir sua identificação por profissional especializado.

\textbf{6. RISCOS}

O principal risco desta pesquisa é a possibilidade, remota, de quebra de sigilo
de seus dados. Para minimizar este risco, adotamos: (a) criptografia dos dados;
(b) acesso restrito por senha; (c) auditoria de acessos; (d) notificação imediata
à Autoridade Nacional de Proteção de Dados (ANPD) e ao(a) senhor(a) em caso de
incidente de segurança.

\textbf{7. BENEFÍCIOS}

O(a) senhor(a) não terá benefício direto. Sua participação contribuirá para o
avanço da ciência odontológica e para o desenvolvimento de ferramentas que
poderão beneficiar futuros pacientes.

\textbf{8. DIREITOS DO TITULAR DOS DADOS (LGPD, Art. 18)}

A qualquer momento, o(a) senhor(a) poderá: (a) confirmar a existência do
tratamento de seus dados; (b) acessar seus dados; (c) corrigir dados incompletos;
(d) solicitar a eliminação de seus dados (revogação do consentimento);
(e) solicitar a portabilidade de seus dados; (f) revogar seu consentimento a
qualquer momento, sem qualquer prejuízo ao seu atendimento odontológico presente
ou futuro.

\textbf{Para exercer qualquer desses direitos, entre em contato com o
pesquisador responsável} pelo e-mail [e-mail] ou telefone [telefone].

\textbf{9. CUSTOS E RESSARCIMENTO}

Sua participação é gratuita e voluntária. O(a) senhor(a) não terá despesas nem
receberá remuneração.

\textbf{10. CONTATO DO COMITÊ DE ÉTICA}

Comitê de Ética em Pesquisa da [instituição] \\
Endereço: [endereço completo] \\
Telefone: [telefone] | E-mail: [e-mail]

\medskip
\textbf{CONSENTIMENTO (LGPD, Art. 8º)}

Declaro que li (ou me foi lido) e compreendi este Termo de Consentimento.
Compreendo que meus dados de saúde bucal, incluindo minha radiografia
panorâmica, serão utilizados para treinamento de inteligência artificial
conforme descrito acima. Autorizo este uso e declaro ter recebido uma via
idêntica deste documento, assinada pelo pesquisador.

\bigskip
\begin{flushright}
\begin{tabular}{p{5cm}p{5cm}}
  \makebox[4cm]{\hrulefill} & \makebox[4cm]{\hrulefill} \\
  Assinatura do participante & Assinatura do pesquisador \\
  & \\
  \makebox[4cm]{\hrulefill} & \makebox[4cm]{\hrulefill} \\
  Nome completo & Nome completo \\
  & \\
  \makebox[4cm]{\hrulefill} & \makebox[4cm]{\hrulefill} \\
  Data & Data \\
\end{tabular}
\end{flushright}
\end{quote}

\teste{O TCLE acima atende a 100\% dos requisitos da LGPD (Art. 8º, 9º e 18),
da Resolução CNS 466/2012 (itens II.23 e IV.3) e da Resolução CNS 510/2016
(Art. 17). Inclui a base legal do tratamento (consentimento), a finalidade
específica, os dados coletados, o prazo de armazenamento, os direitos do titular
e o contato do CEP.}

%-----------------------------------------------------------------------------
% SEÇÃO 25.10
%-----------------------------------------------------------------------------
\section{Síntese do Capítulo}

Este capítulo percorreu o complexo terreno da ética, da proteção de dados e da
regulação da IA em Odontologia. Os principais pontos a reter são:

\begin{itemize}
  \item A LGPD classifica dados odontológicos -- radiografias, imagens
    intraorais, modelos 3D, prontuários -- como dados pessoais sensíveis. Seu
    tratamento exige consentimento explícito ou enquadramento em bases legais
    específicas (tutela da saúde, pesquisa anonimizada).

  \item A anonimização de imagens odontológicas é tecnicamente desafiadora e,
    em muitos casos, a pseudonimização com controle de acesso é a estratégia
    mais realista.

  \item O viés algorítmico é um problema real e documentado: 87\% dos
    \textit{datasets} de IA odontológica não reportam a composição étnico-racial.
    Modelos treinados em populações europeias ou asiáticas podem ter desempenho
    inferior em pacientes brasileiros -- particularmente negros, indígenas e
    pardos.

  \item O PL 2338/2023 classifica sistemas de IA diagnóstica como alto risco e
    estabelece obrigações de transparência, supervisão humana, avaliação de
    impacto e não-discriminação que impactarão diretamente o consultório
    odontológico.

  \item FDA (510k) e CE Mark (MDR + AI Act) são os caminhos regulatórios para
    \textit{software} odontológico com IA nos mercados norte-americano e europeu.

  \item A responsabilidade civil por danos decorrentes de decisões baseadas em
    IA recai, primariamente, sobre o cirurgião-dentista. A IA é ferramenta
    auxiliar, não substituta do julgamento profissional.

  \item A Declaração de Helsinki e as resoluções CNS 466/2012 e 510/2016
    estabelecem o arcabouço ético para pesquisa com IA em Odontologia no Brasil.
\end{itemize}

\destaque{A conformidade ética e regulatória não é um obstáculo à inovação -- é
a condição para que a inovação seja sustentável, confiável e socialmente
legítima. O OdontoIA existe para demonstrar que é possível fazer IA odontológica
de ponta sem jamais transigir com a ética ou com a lei.}

No próximo capítulo, exploraremos como a IA está transformando a educação
odontológica -- dos simuladores inteligentes à gamificação, do Odontinho® como
ferramenta educacional ao futuro do currículo de graduação.

\endinput
